\documentclass[../../main.tex]{subfiles}

\begin{document}

\section{习题}\label{section:ch11-习题}

\begin{exercise}\label{exercise:ch11-1}
证明 $p$-系统的如下推广是严格双曲组，若 $p',q'>0$：
\begin{align*}
\begin{cases}
u_{1,t}-q(u_2)_x=0,\\
u_{2,t}-p(u_1)_x=0.
\end{cases}
\end{align*}
\end{exercise}

\begin{exercise}[a]\label{exercise:ch11-2} 验证浅水波方程（\S11.1\cref{example:ch11-浅水波方程}）构成严格双曲组，只要 $h>0$。

(b) 证明对光滑解 $(h,q)=(h,vh)$，$h>0$，浅水波方程可改写为如下另一种守恒律形式：
\begin{align}
\begin{cases}
h_t+(vh)_x=0,\\
v_t+\left(\dfrac{v^2}{2}+h\right)_x=0.
\end{cases}
\label{align:ch11-浅水交替形式}
\end{align}
验证这是严格双曲组。
\end{exercise}

\begin{exercise}\label{exercise:ch11-3}
对 $z\in\mathbb{R}$，$z\neq 0$，定义矩阵函数
\begin{align*}
\mathbf{B}(z)\coloneqq e^{-1/z^2}\begin{pmatrix}\cos(2/z) & \sin(2/z)\\ \sin(2/z) & -\cos(2/z)\end{pmatrix},
\end{align*}
并设 $\mathbf{B}(0)=0$。证明 $\mathbf{B}$ 是 $C^\infty$ 的且具有实特征值，但我们找不到在 $z=0$ 附近连续依赖于 $z$ 的单位长度右特征向量 $\{\mathbf{r}_1(z),\mathbf{r}_2(z)\}$。当 $z\to 0$ 时特征空间发生了什么？
\end{exercise}

\begin{exercise}[稀疏波曲线。]\label{exercise:ch11-4}证明对于问题 2 中形式 $(*)$ 的浅水波方程，$(h,v)$ 平面中的稀疏波曲线 $R_1,R_2$ 由公式
\begin{align*}
2\sqrt{h}\pm v=c
\end{align*}
给出，其中 $c$ 为常数。
\end{exercise}

\begin{exercise}[激波曲线。]\label{exercise:ch11-5}对于变量 $(h,q)$ 中的浅水波方程，求描述 $(h,q)$ 平面中激波集 $S(z_0)$ 的公式，其中 $z_0=(h_0,q_0)$，$h_0>0$。对于 $q_0=0$ 的特殊情形，证明激波集由表达式
\begin{align*}
q=\pm(h-h_0)h\sqrt{2}\left(\frac{1}{h}+\frac{1}{h_0}\right)^{1/2}
\end{align*}
给出。
\end{exercise}

\begin{exercise}[续上题。]\label{exercise:ch11-6}构造浅水波方程的熵解 $\mathbf{u}=(h,q)$，对应于 Riemann 初值 $u_l=(h_l,0)$（$x<0$）和 $u_r=(h_r,0)$（$x>0$），其中 $h_l>h_r>0$。

你的答案应包括连接 $u_l$ 到中间状态 $u_i=(h_i,q_i)=(h_i,v_i h_i)$ 的稀疏波，以及连接 $u_i$ 到 $u_r$ 的激波。为此将问题 5 的公式转移到 $(h,v)$ 平面，并证明存在过 $(h_l,0)$ 的稀疏波曲线与过 $(h_r,0)$ 的激波曲线的交点 $(h_i,v_i)$。
\end{exercise}

\begin{exercise}\label{exercise:ch11-7}
确认 \S11.3.1 中为正压气体动力学计算的函数 $w_1,w_2$ 确实是 Riemann 不变量。
\end{exercise}

\begin{exercise}\label{exercise:ch11-8}
假设 $\Phi$ 是问题 2 中形式 $(*)$ 的浅水波方程的熵。证明
\begin{align*}
\frac{\partial^2\Phi}{\partial v^2}=h\frac{\partial^2\Phi}{\partial h^2}.
\end{align*}
\end{exercise}

\begin{exercise}\label{exercise:ch11-9}
证明 $\Phi=\rho v^2/2+P(\rho)$ 是正压 Euler 方程（\S11.3.1）的熵，前提是 $P''(\rho)=p'(\rho)/\rho$，$\rho>0$。确认 $\Phi$ 在适当变量中是凸的。相应的熵通量 $\Psi$ 是什么？
\end{exercise}

\begin{exercise}[梯度通量函数。]\label{exercise:ch11-10}假设 $\mathbf{F}=D\phi$，其中 $\phi:\mathbb{R}^m\to\mathbb{R}$。证明
\begin{align*}
\Phi\coloneqq\frac{1}{2}|z|^2,\quad \Psi\coloneqq z\cdot D\phi(z)-\phi(z)
\end{align*}
构成守恒律组 $\mathbf{u}_t+\mathbf{F}(\mathbf{u})_x=0$ 的熵/熵通量对。
\end{exercise}

\begin{exercise}[反梯度通量函数。]\label{exercise:ch11-11}
(a) 假设 $m=2$ 且 $\mathbf{F}$ 具有``反梯度''形式 $\mathbf{F}=(\phi_{z_2},\phi_{z_1})$，其中 $\phi:\mathbb{R}^2\to\mathbb{R}$ 凸。确认
\begin{align*}
\Phi\coloneqq\phi,\quad \Psi\coloneqq\phi_{z_1}\phi_{z_2}
\end{align*}
是方程组 $\mathbf{u}_t+\mathbf{F}(\mathbf{u})_x=0$ 的熵/熵通量对。

(b) 为问题 1 中 $p$-系统的推广找到熵/熵通量对。
\end{exercise}

\begin{exercise}\label{exercise:ch11-12}
阐述 $\Phi,\boldsymbol{\Psi}$ 是一般守恒律组
\begin{align*}
\mathbf{u}_t+\operatorname{div}\mathbf{F}(\mathbf{u})=0 \quad \text{在 } \mathbb{R}^n\times(0,\infty) \text{ 中}
\end{align*}
的熵/熵通量对意味着什么。
\end{exercise}

\begin{exercise}\label{exercise:ch11-13}
非线性介质的 Maxwell 方程部分写作
\begin{align*}
\begin{cases}
\mathbf{D}_t=\operatorname{curl}\mathbf{H},\\
\mathbf{B}_t=-\operatorname{curl}\mathbf{E},
\end{cases}
\end{align*}
在 $\mathbb{R}^3\times(0,\infty)$ 中，其中 $\mathbf{E}$ 为电场，$\mathbf{H}$ 为磁场，$\mathbf{D}$ 为电位移，$\mathbf{B}$ 为磁感应。（假设电流 $\mathbf{J}$ 为零。）我们引入本构关系
\begin{align*}
\mathbf{H}=D_{\mathbf{B}}\eta(\mathbf{B},\mathbf{D}),\quad \mathbf{E}=D_{\mathbf{D}}\eta(\mathbf{B},\mathbf{D}),
\end{align*}
由给定的凸电磁场能量 $\eta=\eta(\mathbf{B},\mathbf{D})$ 确定。证明
\begin{align*}
\Phi\coloneqq\eta,\quad \boldsymbol{\Psi}\coloneqq D_{\mathbf{D}}\eta\times D_{\mathbf{B}}\eta
\end{align*}
是熵/熵通量对。
\end{exercise}

\begin{exercise}[非凸通量函数。]\label{exercise:ch11-14}假设 $u$ 是标量守恒律 $u_t+F(u)_x=0$ 的熵解，且如 \S3.4.1 中，$u$ 在曲线 $\{x=s(t)\}$ 的两侧光滑。

(a) 证明沿此曲线 $u$ 的左、右极限满足
\begin{align*}
F(\lambda u_l+(1-\lambda)u_r)\geqslant\lambda F(u_l)+(1-\lambda)F(u_r) \quad \text{若 } u_l\leqslant u_r
\end{align*}
和
\begin{align*}
F(\lambda u_l+(1-\lambda)u_r)\leqslant\lambda F(u_l)+(1-\lambda)F(u_r) \quad \text{若 } u_r\leqslant u_l
\end{align*}
对每个 $0\leqslant\lambda\leqslant 1$。这些不等式称为\textbf{条件 E}。画图说明这些不等式的几何意义。（与\cref{exercise:ch10-间断面条件} 比较。）

(b) 若 $F$ 一致凸，条件 E 蕴含什么？
\end{exercise}

\begin{exercise}[非凸通量的显式解。]\label{exercise:ch11-15}计算 Riemann 问题
\begin{align}
\begin{cases}
u_t+(u^3-u)_x=0 & \text{在 } \mathbb{R}\times(0,\infty) \text{ 中},\\
u=g & \text{在 } \mathbb{R}\times\{t=0\} \text{ 上},
\end{cases}
\label{align:ch11-非凸通量Riemann}
\end{align}
的唯一熵解，其中
\begin{align*}
g(x)=\begin{cases}-1 & \text{若 } x<0,\\ 1 & \text{若 } x>0.\end{cases}
\end{align*}
（提示：解涉及沿左边缘带激波的稀疏波。确保你的解验证了上题引入的条件 E。）
\end{exercise}

\begin{exercise}[续上题。]\label{exercise:ch11-16}讨论问题 15 中守恒律 $(*)$ 的熵解结构，其中
\begin{align*}
g(x)=\begin{cases}
-1+\lambda & \text{若 } x<-1,\\
-1-\lambda x & \text{若 } -1<x<0,\\
1 & \text{若 } x>0,
\end{cases}
\end{align*}
当 $\lambda>0$ 小时。此扰动如何影响上题所求解中激波的行为？
\end{exercise}

\begin{exercise}\label{exercise:ch11-17}
假设 $u$ 是多变量标量守恒律的光滑解，具有空间紧支集：
\begin{align*}
\begin{cases}
u_t+\operatorname{div}\mathbf{F}(u)=0 & \text{在 } \mathbb{R}^n\times(0,\infty) \text{ 中},\\
u=g & \text{在 } \mathbb{R}^n\times\{t=0\} \text{ 上},
\end{cases}
\end{align*}
其中 $\mathbf{F}=(F_1,\cdots,F_n)$。证明若 $\widehat{u}$ 是具有初值 $\widehat{g}$ 的另一个光滑解，我们有 $L^1$ 压缩估计
\begin{align*}
\|u(\cdot,t)-\widehat{u}(\cdot,t)\|_{L^1(\mathbb{R}^n)}\leqslant\|g-\widehat{g}\|_{L^1(\mathbb{R}^n)}
\end{align*}
对每个时间 $t\geqslant 0$。

（提示：证明 $w\coloneqq u-\widehat{u}$ 满足 $w_t+\operatorname{div}(\mathbf{b}w)=0$ 对适当向量场 $\mathbf{b}=\mathbf{b}(x,t)$。令 $T>0$ 并引入伴随问题 $v_t+\mathbf{b}\cdot Dv=0$ 在 $\mathbb{R}^n\times[0,T]$ 中的解 $v$，终端条件 $v(\cdot,T)=\operatorname{sgn}(w(\cdot,T))$。）
\end{exercise}

\end{document}